\subsection{Guardrail Semantics}\label{subsection:semantics}

The previous subsections introduced guardrails as an abstraction and GDL as a language for specifying them. We now give a simple semantics for guardrails that captures their effect on policy execution. The purpose of this semantics is to formalize \emph{what} a guardrail does once it is evaluated. To keep the presentation as simple as possible, we do not model performance-related aspects of guardrail evaluation, such as trigger conditions, as they do not change the logical effect of guardrail evaluation itself.

We consider the application of a guardrail to a policy $\pi$ in a system $S$, where the system may be the OS as a whole or a particular subsystem in which the policy operates. Let $\Sigma$ denote the set of possible run-time states of the system, including the relevant kernel-visible statistics and the state of in-kernel objects. We write $\Pi$ to denote the set of policies that may be deployed in the system. 
%A deployed policy may additionally carry parameters, but we leave these implicit unless they are needed explicitly.

To describe the semantics, a guardrail is modeled as a triple
$ 
G = (P, A, E)
$ 
where:
$P : \Sigma \to \{0,1\}$ is the run-time condition monitored by the guardrail,  $A : \Sigma \times \Pi \to \Pi$ is the corrective action applied when the condition fails, and 
$E : \Pi \to \{0,1\}$ specifies whether the guardrail is enabled for a given policy. Intuitively, $P$ captures the condition that should continue to hold while the policy is operating in its intended regime, $A$ specifies how the policy should be repaired when that condition is violated, and $E$ determines the policies to which the guardrail applies.

We write
$ 
G \models_{\sigma} \pi \to \pi'
$
to mean that guardrail $G$, when evaluated in state $\sigma$ under deployed policy $\pi$, produces an updated policy $\pi'$. This relation is defined by the following rules:
\begin{mathpar}
\inferrule[Guard-Disabled]{
    \enabled(\policy_\args) = 0
}{
  \guard{(\property{},\action{},\enabled)}{\sysstate}{\ppol{\policy}{\args}}{\ppol{\policy}{\args}}
}\quad
\inferrule[Guard-OK]{ 
    \enabled(\policy_\args) = 1 =
    \property{}(\sysstate)
}{
  \guard{(\property{},\action{},\enabled)}{\sysstate}{\ppol{\policy}{\args}}{\ppol{\policy}{\args}}
}\\
\inferrule[Guard-Violation]{
    \enabled(\policy_\args)
    \neq 0 = \property{}(\sysstate)
}{
    \guard{(\property{},\action{},\enabled)}{\sysstate}{\ppol{\policy}{\args}}{\action{}(\sysstate,\ppol{\policy}{\args})}
}
\end{mathpar}

That is,  a disabled guardrail has no effect, an enabled guardrail leaves the policy unchanged when its condition already holds, and otherwise applies its corrective action.


\paragraph{Guarded-system semantics.}
We now define the overall system semantics induced by a single guardrail $G$.
Let a system be given by the tuple $(\states,\policies,\transitions)$, where $\states$ and $\policies$ are defined as above, and $\transitions$ is the transition relation for the system. That is, $\TransitionsVia{\sysstate}{\ppol\policy\args}{\sysstate'}$ indicates that the system steps from $\sysstate$ to $\sysstate'$ under policy $\ppol\policy\args$.
The guarded system is then an extension of $S$, where at each step, it first checks whether the guardrail is violated and repairs the deployed parameterized policy if needed.
This is captured by the rule:
\begin{mathpar}
\inferrule[Guarded-Transition]{
    \TransitionsVia{\sysstate}{\ppol{\policy}{\args'}}{\sysstate'}
    \quad
    \guard{G}{{\sysstate'}}{\ppol{\policy}{\args}}{\ppol{\policy'}{\args'}} 
}{
    \transition{G}{\runtimestate{\sysstate}{\ppol{\policy}{\args}}}{\runtimestate{\sysstate'}{\ppol{\policy'}{\args'}}}
}
\end{mathpar}

This semantics captures the core logical role of a guardrail: it interposes between the current system state and the next policy-controlled transition, possibly changing the deployed policy before that transition occurs. Multiple guardrails can be incorporated  by composing this interposition step in some  nondeterministic order.

\subsection{Guardrail Semantics}\label{subsection:semantics}

The previous subsections introduced guardrails as an abstraction and GDL as a language for specifying them. We now give a simple semantics for guardrails that captures their effect on policy execution. The purpose of this semantics is to formalize \emph{what} a guardrail does once it is evaluated, and then use these semantics to build the verification framework in \Cref{sec:verification}.
% To keep the presentation focused, we abstract away from performance-related aspects of evaluation, such as trigger timing, and model only the logical effect of evaluating a guardrail on a state and a deployed policy.

We consider the application of a guardrail to a policy $\pi$ in a system $S$, where the system may be the OS as a whole or a particular subsystem in which the policy operates. Let $\Sigma$ denote the set of possible run-time states of the system, including the relevant kernel-visible statistics and the state of in-kernel objects, and let $\mathsf P$ denote the set of policy names and $V^*$ denote the set of (positional) argument lists. The set of policies is $\Pi = \mathsf P \times V^*$.

A guardrail is modeled as a triple
$
G = (P,A,E),
$
where $P : \Sigma \to \{0,1\}$ is the run-time condition monitored by the guardrail, $A : \Sigma \times \Pi \to \Pi$ is the corrective action applied when that condition fails, and $E : \Pi \to \{0,1\}$ specifies whether the guardrail is enabled for a given policy. Intuitively, $P$ captures the condition that should continue to hold while the policy is operating in its intended regime, $A$ specifies how the policy is repaired when that condition is violated, and $E$ determines the policies to which the guardrail applies.

We write
$ 
G \models_{\sigma} \pi \to \pi'
$
to mean that guardrail $G$, when evaluated in state $\sigma$ under deployed policy $\pi$, produces an updated policy $\pi'$. This relation is defined by the  rules:
\begin{mathpar}
\inferrule[Policy-Same]{
    E(\pi)=0 \lor P(\sigma)
}{
    G \models_{\sigma} \pi \to \pi
}
\and
\inferrule[Policy-Change]{
    E(\pi)=1 \quad \neg P(\sigma)
}{
    G \models_{\sigma} \pi \to A(\sigma,\pi)
}
\end{mathpar}

Thus, the policy remains unchanged if a guardrail is disabled or a guardrail is enabled but its condition already holds; otherwise, the corrective action changes the policy.

A guardrail $G = (E,A,P)$ monitors a system-under-observation given by a transition relation $\TransitionsVia{\sysstate}{\policy}{\sysstate'}$, which indicates that under policy $\policy$, the system steps from $\sysstate$ to $\sysstate'$. A single guardrail composes with the system semantics as follows:
\[
\inferrule*[right={~~Guarded-Trans}]
  { \sigma \xrightarrow{\pi} \sigma' \\
    G \models_{\sigma'} \pi \to \pi' }
  { G \models \langle \sigma,\pi\rangle  \rightsquigarrow \langle \sigma',\pi'\rangle }
\]
In other words, the system takes a step according to policy $\policy$, the guardrail inspects the new state $\sysstate$, and updates the policy to $\policy'$, which may be a different policy (via \textsc{Policy-Change}) or the same policy (via \textsc{Policy-Same}).

Our system can also support a set of guardrails monitoring a system simultaneously. The formal semantics for this can be found in Appendix~\ref{appdx:full-semantics}.